\documentclass[11pt]{article}

\usepackage[margin=1in]{geometry}
\usepackage{amsmath,amssymb}
\usepackage{enumitem}
\usepackage[hidelinks]{hyperref}

\setlist[itemize]{leftmargin=1.4em}
\setlist[enumerate]{leftmargin=1.6em}

\title{SYSC 5108 Exam Study\\Assignment 4, Question 4}
\author{Jesse Levine}
\date{}

\begin{document}
\maketitle

\section*{Question}

In the grid world example, rewards are positive for the gold state, negative for
traps, and zero for all other states. Are the absolute values of these rewards
important, or only the intervals between them? Prove that adding a constant
\(c\) to all rewards adds a constant, \(v_c\), to the values of all states, and
thus does not affect the relative values of any states under any policy. What is
\(v_c\) in terms of \(c\) and \(\gamma\)?

\section*{Course and Textbook Cross-References}

\begin{itemize}
    \item \textbf{Chapter 6 slides, Reward and Return:} the discounted return is
    defined as
    \[
        G_t = R_{t+1} + \gamma R_{t+2} + \gamma^2 R_{t+3} + \cdots.
    \]
    \item \textbf{Chapter 6 slides, Value Function and Bellman Equation:} a
    state value is the expected discounted return under a policy.
    \item \textbf{Chapter 6 slides, Grid World example:} the state ranking comes
    from comparing values, not from the raw reward labels alone.
    \item \textbf{Goodfellow, Bengio, and Courville, \emph{Deep Learning},
    reinforcement-learning discussion:} a policy is evaluated through expected
    reward/return, so changing every reward by the same additive offset shifts the
    return baseline uniformly.
\end{itemize}

\section*{What the Question Is Really Asking}

The question is about whether a \emph{uniform reward shift}
\[
    R'_{t+1} = R_{t+1} + c
\]
changes which states are better or worse under a fixed policy \(\pi\).

The answer is:
\begin{itemize}
    \item the \emph{absolute numerical labels} of the rewards are not what matter,
    \item what matters for comparing states is their \emph{relative difference in
    value},
    \item adding the same constant \(c\) to every reward shifts every state value
    by the same constant,
    \item so the ordering of states does not change.
\end{itemize}

\section*{Step 1: Start from the Definition of Value}

Under policy \(\pi\), the state-value function is
\[
    v_\pi(s) = \mathbb{E}_\pi[G_t \mid S_t = s],
\]
where the discounted return is
\[
    G_t = \sum_{k=0}^{\infty} \gamma^k R_{t+k+1}.
\]

Now define the shifted rewards:
\[
    R'_{t+1} = R_{t+1} + c.
\]
Then the new return is
\[
    G'_t = \sum_{k=0}^{\infty} \gamma^k R'_{t+k+1}.
\]

\section*{Step 2: Expand the New Return}

Substitute \(R'_{t+k+1} = R_{t+k+1} + c\):
\[
\begin{aligned}
    G'_t
    &= \sum_{k=0}^{\infty} \gamma^k (R_{t+k+1} + c) \\
    &= \sum_{k=0}^{\infty} \gamma^k R_{t+k+1}
     + \sum_{k=0}^{\infty} \gamma^k c \\
    &= G_t + c \sum_{k=0}^{\infty} \gamma^k.
\end{aligned}
\]

For \(0 \le \gamma < 1\), the geometric series is
\[
    \sum_{k=0}^{\infty} \gamma^k = \frac{1}{1-\gamma}.
\]
Therefore,
\[
    G'_t = G_t + \frac{c}{1-\gamma}.
\]

\section*{Step 3: Convert the Return Shift into a Value Shift}

Take expectation under the same policy \(\pi\):
\[
\begin{aligned}
    v'_\pi(s)
    &= \mathbb{E}_\pi[G'_t \mid S_t = s] \\
    &= \mathbb{E}_\pi\left[G_t + \frac{c}{1-\gamma}\,\middle|\,S_t=s\right] \\
    &= \mathbb{E}_\pi[G_t \mid S_t=s] + \frac{c}{1-\gamma} \\
    &= v_\pi(s) + \frac{c}{1-\gamma}.
\end{aligned}
\]

So the constant added to every state value is
\[
    \boxed{v_c = \frac{c}{1-\gamma}}.
\]

\section*{Step 4: Show that Relative Values Do Not Change}

Take any two states \(s_1\) and \(s_2\). Under the shifted rewards,
\[
\begin{aligned}
    v'_\pi(s_1) - v'_\pi(s_2)
    &=
    \left(v_\pi(s_1) + \frac{c}{1-\gamma}\right)
    -
    \left(v_\pi(s_2) + \frac{c}{1-\gamma}\right) \\
    &= v_\pi(s_1) - v_\pi(s_2).
\end{aligned}
\]

So the difference between any pair of state values is unchanged. That means:
\begin{itemize}
    \item if \(v_\pi(s_1) > v_\pi(s_2)\) before the shift, then
    \(v'_\pi(s_1) > v'_\pi(s_2)\) after the shift,
    \item if two states had the same value before, they still have the same value
    after,
    \item the policy-induced ranking of states is unchanged.
\end{itemize}

\section*{Interpretation for Grid World}

In the grid world, suppose gold has positive reward, traps have negative reward,
and ordinary states have zero reward. If we add the same constant \(c\) to
\emph{every} reward:
\begin{itemize}
    \item gold becomes ``old gold reward \(+c\),''
    \item traps become ``old trap reward \(+c\),''
    \item ordinary states become ``\(0+c\).''
\end{itemize}

But every trajectory receives the same extra discounted baseline
\[
    \frac{c}{1-\gamma},
\]
so every state's value shifts by exactly that same amount. Therefore the agent's
preference between states does not change.

\section*{Final Answer}

Only the \emph{intervals between rewards} matter for comparing state values under
a fixed policy, not their absolute baseline. If every reward is shifted by the
same constant \(c\), then for any state \(s\),
\[
    v'_\pi(s) = v_\pi(s) + v_c,
\]
where
\[
    \boxed{v_c = \frac{c}{1-\gamma}} \qquad (0 \le \gamma < 1).
\]
Because the same constant is added to all states,
\[
    v'_\pi(s_1) - v'_\pi(s_2) = v_\pi(s_1) - v_\pi(s_2),
\]
so the relative values of states are unchanged under any policy.

\section*{Exam Shortcut}

If a question asks what happens when all rewards are replaced by \(R+c\), write:
\[
    G'_t = G_t + c(1+\gamma+\gamma^2+\cdots)
         = G_t + \frac{c}{1-\gamma}.
\]
Then immediately conclude:
\[
    v'_\pi(s) = v_\pi(s) + \frac{c}{1-\gamma},
\]
so \textbf{state ordering does not change}.

\section*{Sources Used}

\begin{itemize}
    \item Assignment 4 PDF, Question 4.
    \item Local submitted Assignment 4 PDF checked for consistency of the core
    result.
    \item Chapter 6 slides on reward/return, value functions, Bellman equations,
    and the grid world example.
    \item Goodfellow, Bengio, and Courville, \emph{Deep Learning}, reinforcement-learning discussion on policies, rewards, and return-based evaluation.
\end{itemize}

\end{document}
